% Options for packages loaded elsewhere
\PassOptionsToPackage{unicode}{hyperref}
\PassOptionsToPackage{hyphens}{url}
%
\documentclass[
  11pt,
]{article}
\usepackage{amsmath,amssymb}
\usepackage{iftex}
\ifPDFTeX
  \usepackage[T1]{fontenc}
  \usepackage[utf8]{inputenc}
  \usepackage{textcomp} % provide euro and other symbols
\else % if luatex or xetex
  \usepackage{unicode-math} % this also loads fontspec
  \defaultfontfeatures{Scale=MatchLowercase}
  \defaultfontfeatures[\rmfamily]{Ligatures=TeX,Scale=1}
\fi
\usepackage{lmodern}
\ifPDFTeX\else
  % xetex/luatex font selection
\fi
% Use upquote if available, for straight quotes in verbatim environments
\IfFileExists{upquote.sty}{\usepackage{upquote}}{}
\IfFileExists{microtype.sty}{% use microtype if available
  \usepackage[]{microtype}
  \UseMicrotypeSet[protrusion]{basicmath} % disable protrusion for tt fonts
}{}
\makeatletter
\@ifundefined{KOMAClassName}{% if non-KOMA class
  \IfFileExists{parskip.sty}{%
    \usepackage{parskip}
  }{% else
    \setlength{\parindent}{0pt}
    \setlength{\parskip}{6pt plus 2pt minus 1pt}}
}{% if KOMA class
  \KOMAoptions{parskip=half}}
\makeatother
\usepackage{xcolor}
\usepackage[margin=0.82in]{geometry}
\usepackage{longtable,booktabs,array}
\usepackage{calc} % for calculating minipage widths
% Correct order of tables after \paragraph or \subparagraph
\usepackage{etoolbox}
\makeatletter
\patchcmd\longtable{\par}{\if@noskipsec\mbox{}\fi\par}{}{}
\makeatother
% Allow footnotes in longtable head/foot
\IfFileExists{footnotehyper.sty}{\usepackage{footnotehyper}}{\usepackage{footnote}}
\makesavenoteenv{longtable}
\usepackage{graphicx}
\makeatletter
\def\maxwidth{\ifdim\Gin@nat@width>\linewidth\linewidth\else\Gin@nat@width\fi}
\def\maxheight{\ifdim\Gin@nat@height>\textheight\textheight\else\Gin@nat@height\fi}
\makeatother
% Scale images if necessary, so that they will not overflow the page
% margins by default, and it is still possible to overwrite the defaults
% using explicit options in \includegraphics[width, height, ...]{}
\setkeys{Gin}{width=\maxwidth,height=\maxheight,keepaspectratio}
% Set default figure placement to htbp
\makeatletter
\def\fps@figure{htbp}
\makeatother
\setlength{\emergencystretch}{3em} % prevent overfull lines
\providecommand{\tightlist}{%
  \setlength{\itemsep}{0pt}\setlength{\parskip}{0pt}}
\setcounter{secnumdepth}{-\maxdimen} % remove section numbering
\ifLuaTeX
\usepackage[bidi=basic]{babel}
\else
\usepackage[bidi=default]{babel}
\fi
\babelprovide[main,import]{american}
% get rid of language-specific shorthands (see #6817):
\let\LanguageShortHands\languageshorthands
\def\languageshorthands#1{}
\usepackage{microtype}
\usepackage{etoolbox}
\usepackage{xcolor}
\usepackage{needspace}

\definecolor{guideblue}{HTML}{334E68}
\definecolor{guideteal}{HTML}{2C7A7B}

\setlength{\parindent}{0pt}
\setlength{\parskip}{0.55em plus 0.12em minus 0.08em}
\setlength{\emergencystretch}{3em}
\renewcommand{\arraystretch}{1.13}
\AtBeginEnvironment{longtable}{\small}

\AtBeginDocument{%
  \hypersetup{
    colorlinks=true,
    linkcolor=guideblue,
    citecolor=guideteal,
    urlcolor=guideteal,
    pdftitle={Instrumental Variables Before and After an Event: Identification with Endogenous Exposure},
    pdfauthor={Sriteja Burla}
  }%
}
\ifLuaTeX
  \usepackage{selnolig}  % disable illegal ligatures
\fi
\IfFileExists{bookmark.sty}{\usepackage{bookmark}}{\usepackage{hyperref}}
\IfFileExists{xurl.sty}{\usepackage{xurl}}{} % add URL line breaks if available
\urlstyle{same}
\hypersetup{
  pdftitle={Instrumental Variables Before and After an Event},
  pdfauthor={Sriteja Burla},
  pdflang={en-US},
  hidelinks,
  pdfcreator={LaTeX via pandoc}}

\title{Instrumental Variables Before and After an Event}
\usepackage{etoolbox}
\makeatletter
\providecommand{\subtitle}[1]{% add subtitle to \maketitle
  \apptocmd{\@title}{\par {\large #1 \par}}{}{}
}
\makeatother
\subtitle{Identification with Endogenous Exposure}
\author{Sriteja Burla}
\date{August 2026}

\begin{document}
\maketitle

{
\setcounter{tocdepth}{2}
\tableofcontents
}
\hypertarget{summary}{%
\subsection{Summary}\label{summary}}

Suppose a hospital is acquired. Did the acquisition change the effect of
receiving care there? A direct comparison with patients treated
elsewhere is difficult because hospital destination is not random.
Ambulance dispatch offers a possible instrument: otherwise similar
patients may be assigned to companies whose hospital preferences shift
which hospital initially receives them (Doyle et al.~2015).

One could estimate the ambulance IV before and after acquisition and
subtract the Wald ratios. Each ratio can identify a local effect for
patients whose destination responds to the instrument in that period. To
attribute the difference to acquisition, the estimates must refer to the
same target population and comparable compliers. We also need to specify
how the LATE would have evolved without acquisition, with exclusion
holding in that counterfactual state as well as in the observed periods.

The benchmark in this guide assumes that the same units comply before
the event and under both post-event states, and that the LATE would have
remained at its pre-event value without the event. Under the full
conditions stated below, subtracting the standardized pre-event Wald
ratio from the standardized post-event ratio then identifies the event
effect. Sensitivity analysis allows departures from the assumptions
about who complies and what the counterfactual LATE would have been.

\hypertarget{the-empirical-setting}{%
\subsection{1. The empirical setting}\label{the-empirical-setting}}

\hypertarget{roles-of-the-event-and-instrument}{%
\subsubsection{1.1. Roles of the event and
instrument}\label{roles-of-the-event-and-instrument}}

I use a hospital acquired at a known date as the running example: the
acquisition is the event, and emergency patients are the units whose
hospital exposure may change. Let \(D=1\) denote initial care at the
focal hospital, while \(D=0\) denotes a prespecified alternative
exposure---for example, initial care at one comparison hospital or at
hospitals drawn from a fixed eligible-hospital distribution. Let
\(Z\in\{0,1\}\) denote a prespecified encouragement derived from
ambulance-company assignment or company preferences.

Ambulance assignment shifts hospital destination within a period.
Acquisition may change care at the focal hospital and at the
alternatives used for comparison. The role of \(Z\) is therefore to
identify the effect of hospital exposure in each period. Attributing a
change in that effect to acquisition requires assumptions that link the
periods.

In fuzzy or instrumented DiD, the instrument comes from a policy or
group change (de Chaisemartin and D'Haultfœuille 2018; Miyaji 2024).
Here, the instrument operates within each period. The event changes the
environment in which the exposure effects are compared.

\Needspace{0.62\textheight}

\textbf{Figure 1. The event and the instrument play different causal
roles}

\includegraphics[width=0.88\textwidth,height=\textheight]{figures/figure1_nested_event_roles.pdf}

\emph{Notes:} The IV shifts exposure within a period. The event may
change both the exposure process and the effect of exposure. Exclusion
rules out an effect of \(Z\) on \(Y\) through channels outside the
prespecified exposure contrast.

The same problem arises when a workplace policy prompts workers to sort
across establishments, when families choose among schools that adopt
programs, or when a reform changes how courts assign cases.

\hypertarget{target-population-and-intervention-definitions}{%
\subsubsection{1.2. Target population and intervention
definitions}\label{target-population-and-intervention-definitions}}

Define the \textbf{target population} by a rule that remains fixed
across the event and throughout the analyzed periods. In the hospital
application, it might include all eligible emergency patients
originating in prespecified pickup areas around the hospital, whether
they receive care at the focal hospital or are ultimately taken
elsewhere.

Determine eligibility before ambulance assignment. Defining the study
population by where patients actually receive care would condition on
the exposure that the instrument changes. Predetermined market, cohort,
or source-population indicators can refine the target population. Fixed
effects for the provider that actually treats a patient raise a
different problem: they condition on an endogenous response and may
absorb the first stage or select patients on unobserved determinants of
outcomes. Konetzka, Yang, and Werner (2019) discuss the related problem
of using a patient-level instrument to study a provider-level attribute.

The instrument and exposure must mean the same thing in every period.
The two values of \(Z\) should represent the same encouragement, and
\(D=1\) and \(D=0\) should describe the same exposure contrast. If
\(D=0\) pools several hospitals, a change in their shares may change
what alternative care means. Fixing the distribution of alternative
hospitals addresses this problem. When that is infeasible, researchers
must justify treating care at these hospitals as equivalent versions of
the exposure, or distinguish among them in the exposure definition.

The data may be a panel or repeated cross sections; eligibility and the
availability of outcomes may change over time. State how eligible units
enter the records and which outcomes are observed. If exposure affects
survival, attrition, or missingness, it also changes whose outcomes
contribute to the estimate. That selection requires a separate analysis.

\hypertarget{identification-of-the-event-effect}{%
\subsection{2. Identification of the event
effect}\label{identification-of-the-event-effect}}

\hypertarget{period-specific-local-average-treatment-effects}{%
\subsubsection{2.1. Period-specific local average treatment
effects}\label{period-specific-local-average-treatment-effects}}

Let \(t=0\) denote a pre-event period and \(t=1\) a post-event period.
For predetermined covariates \(X=x\), the conditional Wald ratio is

\[
\beta_t(x)=
\frac{E[Y\mid Z=1,X=x,t]-E[Y\mid Z=0,X=x,t]}
{E[D\mid Z=1,X=x,t]-E[D\mid Z=0,X=x,t]}.
\tag{1}
\]

Write the numerator as \(\rho_t(x)\) and the denominator as
\(\pi_t(x)\). Under the usual binary-IV conditions---independence,
exclusion, monotonicity, positive probability for both instrument
values, and relevance---\(\beta_t(x)\) is the exposure LATE for
period-\(t\) compliers whose outcomes are observed (Imbens and Angrist
1994). Define \(Z=1\) as the value intended to increase exposure and
keep that coding in every period.

The IV conditions must hold in the analyzed sample in every period used
in the design. Section 2.2 states the additional requirements when
outcomes are selectively observed or when the estimand covers all
compliers in the target population. Report the instrument propensity,
reduced form, first stage, and evidence on instrument strength for each
period. A pooled first stage can conceal a weak component.

\hypertarget{standardization-and-outcome-observation}{%
\subsubsection{2.2. Standardization and outcome
observation}\label{standardization-and-outcome-observation}}

Raw Wald ratios can change even when the conditional LATEs do not,
simply because the observed covariate distribution of compliers changes.
Choose a reference distribution \(F_X^\star\) before examining outcomes,
then average the conditional ratios over the part of its domain with
common support:

\[
\beta_t^\star=E_{F_X^\star}[\beta_t(X)],
\qquad
\Delta^{IV,\star}=\beta_1^\star-\beta_0^\star.
\tag{2}
\]

The covariate distribution of pre-event compliers is a useful default.
Given the observed pre-event distribution \(F_{X0}\), the reference
measure is proportional to \(\pi_0(x)dF_{X0}(x)\). Restrict it to strata
that support both period-specific ratios. Researchers can choose another
policy-relevant distribution, but should explain that choice and report
any trimming.

Standardization gives the pre- and post-event Wald ratios the same
observed covariate distribution. To extend the interpretation to all
compliers in the target population, their conditional mean exposure
effect must equal that among compliers with observed outcomes. Outcomes
must have a positive probability of being observed in every target
stratum, and exclusion must hold for the target population as well as
for the analyzed records. Selective survival or attrition will often
require an explicit model or sensitivity analysis.

\hypertarget{complier-populations-and-the-post-event-counterfactual}{%
\subsubsection{2.3. Complier populations and the post-event
counterfactual}\label{complier-populations-and-the-post-event-counterfactual}}

Patients with the same observed covariates may still differ in whether
they comply with the instrument. Thus, standardization alone does not
ensure that \(\Delta^{IV,\star}\) compares the same compliers. The
benchmark assumes that the same units would comply before the event and
after it under both the observed and counterfactual states. A weaker
approach allows membership to change if the relevant conditional mean
exposure effects remain equal across the complier populations. Sections
B.2 and B.3 state these conditions.

Even with comparable complier populations, a second question remains:
what would the post-event LATE have been at the same date without the
event? Let \(L_1^{E,\cap,\star}\) denote the standardized post-event
LATE under the event and \(L_1^{0,\cap,\star}\) that missing
counterfactual; both refer to units who comply in all three relevant
states and use the reference covariate distribution. The event estimand
is

\[
\theta^{\cap,\star}=L_1^{E,\cap,\star}-L_1^{0,\cap,\star}.
\tag{3}
\]

After the event, the data reveal only the first term. The simplest
assumption sets the missing counterfactual equal to the comparable
pre-event LATE. With several pre-event estimates, researchers can
instead extrapolate the pre-event LATE path linearly or bound departures
from it. The restriction concerns how the LATE would have evolved. A
parallel-trends assumption about average untreated outcomes does not
supply that restriction.

Exclusion must hold in the counterfactual post-event state too: holding
exposure fixed, \(Z\) has no effect on the outcome that would be
observed without the event. Institutional evidence must support this
restriction because the state is unobserved.

No anticipation gives the pre-event estimates their counterfactual
meaning. Announcements, preparation, or early changes in ambulance
behavior may affect periods close to acquisition. If so, use an earlier
reference period or redefine the transition window.

\Needspace{0.58\textheight}

\textbf{Figure 2. The observed IV path leaves one counterfactual LATE
missing}

\includegraphics[width=0.88\textwidth,height=\textheight]{figures/figure2_missing_late_path.pdf}

\emph{Notes:} The observed pre-event and post-event Wald ratios identify
period-specific local effects under their respective IV conditions. An
event interpretation also requires a comparable complier population and
a specification for how the post-event LATE would have evolved without
the event. The dashed line is schematic: applications may hold the
pre-event LATE fixed, extrapolate its earlier path, or bound plausible
counterfactual values.

Under the benchmark conditions above, including a common complier
population and a counterfactual post-event LATE equal to the pre-event
LATE, the event effect is

\[
\boxed{\theta^{\cap,\star}=\beta_1^\star-\beta_0^\star.}
\tag{4}
\]

A different counterfactual path requires replacing the pre-event ratio
with the corresponding counterfactual value. In either case, the
interpretation also rests on the stated restrictions on outcome
observation, complier membership, anticipation, and exclusion.

\hypertarget{what-the-estimates-identify}{%
\subsubsection{2.4. What the estimates
identify}\label{what-the-estimates-identify}}

Table 1 separates the claims supported at each stage of the design.

\Needspace{0.31\textheight}

\textbf{Table 1. What the design supports at each stage}

\begingroup
\renewcommand{\arraystretch}{1.30}

\begin{longtable}[]{@{}
  >{\raggedright\arraybackslash}p{(\columnwidth - 2\tabcolsep) * \real{0.5000}}
  >{\raggedright\arraybackslash}p{(\columnwidth - 2\tabcolsep) * \real{0.5000}}@{}}
\toprule\noalign{}
\begin{minipage}[b]{\linewidth}\raggedright
Maintained conditions
\end{minipage} & \begin{minipage}[b]{\linewidth}\raggedright
Supported interpretation
\end{minipage} \\
\midrule\noalign{}
\endhead
\bottomrule\noalign{}
\endlastfoot
Outcome and exposure means by instrument value are well defined &
Reduced forms and first stages \\
Period-specific IV conditions hold & A LATE of exposure for each
period's observed compliers \\
Common support and standardization hold, and conditional mean effects
are equal for observed and target-population compliers & Average
conditional exposure effects for target-population compliers,
standardized to the chosen reference distribution \\
The same complier population is used across periods and event states; no
anticipation and exclusion hold; and the counterfactual post-event LATE
is correctly specified & Change in the standardized LATE caused by the
event for that complier population \\
\end{longtable}

\endgroup

When researchers cannot defend the final row, they can still report a
standardized change in period-specific LATEs: exposure effects changed
for the complier populations reached by the instrument. Decide before
estimation which interpretation the design aims to support and what to
report if the stronger assumptions prove unconvincing.

\hypertarget{decomposition-and-sensitivity-analysis}{%
\subsection{3. Decomposition and sensitivity
analysis}\label{decomposition-and-sensitivity-analysis}}

\hypertarget{diagnostic-decomposition}{%
\subsubsection{3.1. Diagnostic
decomposition}\label{diagnostic-decomposition}}

A standardized IV change need not be the event effect: the LATE could
have moved even without the event, the instrument may reach a different
complier population, or the observed records may cover a different part
of the target population. The decomposition separates these
possibilities:

\[
\Delta^{IV,\star}
=\theta^{\cap,\star}
+\Delta^{0,\star}
+B^{C,\star}
+B^{O,\star}.
\tag{5}
\]

Here \(\Delta^{0,\star}\) is the change in the standardized LATE that
would have occurred without the event for the common complier
population. \(B^{C,\star}\) measures how changes in who complies affect
the comparison. \(B^{O,\star}\) measures changes in the gap between
observed compliers and compliers in the target population. The
decomposition presumes that the period-specific IV conditions and
definitions of the interventions are valid.

Historical IV estimates can inform \(\Delta^{0,\star}\) if they use the
same instrument, exposure, target population, and assumptions about
complier comparability. Covariate balance and standardization help
assess changes in observed composition. First-stage changes reveal
changes in the size of the complier group, but equal first stages can
hide substantial entry and exit. Sampling rates, attrition patterns, and
checks on which outcomes are observed help assess \(B^{O,\star}\).

The three gaps can instead be bounded jointly. If their absolute values
are bounded by \(h_0\), \(h_C\), and \(h_O\), then

\[
\theta^{\cap,\star}
\in
\left[\Delta^{IV,\star}-H,\ \Delta^{IV,\star}+H\right],
\qquad H=h_0+h_C+h_O.
\tag{6}
\]

Vary the discrepancies jointly because they may occur together.
Rambachan and Roth (2023) provide a template for calibrating
restrictions on the counterfactual LATE path. Applying that approach
here requires inference for the resulting IV estimand.

\Needspace{0.74\textheight}

\hypertarget{simulation-design-and-results}{%
\subsubsection{3.2. Simulation design and
results}\label{simulation-design-and-results}}

Figure 3 shows why each assumption matters. The simulation follows four
repeated cross sections around one event. The IV is valid in every
period, and the event causes a change in LATE of one in every scenario.
The scenarios separately vary observed covariate composition, unobserved
complier types, and the path the LATE would have followed without the
event.

\Needspace{0.56\textheight}

\textbf{Figure 3. Each assumption addresses a different source of bias}

\includegraphics{figures/figure3_standardization_and_continuation.png}

\emph{Notes:} Markers are means across 2,000 replications; vertical
lines mark the 2.5th and 97.5th percentiles. Each period has 2,500
observations. The dashed line gives the true change in LATE caused by
the event. ``Standardized'' averages over the pre-event complier
covariate distribution. ``Pre-event LATE held fixed'' uses the
comparable pre-event LATE as the counterfactual post-event value;
``Linear extrapolation'' uses three pre-event LATEs. The replication
materials document the data-generating process, fixed seed, and
outputs.\footnote{Replication materials: \url{https://github.com/TejaBurla89/iv_nested_event}.}

In the benchmark case, the same units comply in all three states, and
all three estimators average about one. When only the observed
distribution of \(X\) changes, the unstandardized estimate rises to
1.75; standardizing to the pre-event complier distribution brings it
back to 1.00.

When unobserved complier types turn over, the standardized estimate
remains near 1.80 even though the observed covariate distribution does
not change. In the last scenario, the LATE would have risen without the
event. Holding the pre-event value fixed gives 1.40; linear
extrapolation from three pre-event LATEs gives 0.99.

\hypertarget{multiple-periods-and-staggered-events}{%
\subsection{4. Multiple periods and staggered
events}\label{multiple-periods-and-staggered-events}}

\hypertarget{event-time-local-average-treatment-effects}{%
\subsubsection{4.1. Event-time local average treatment
effects}\label{event-time-local-average-treatment-effects}}

With more periods, estimate a path of standardized LATEs by event time.
Its pre-event portion shows how the local effect was evolving. Report
the first stage and common support alongside these estimates to show how
they changed over time.

The pre-event estimates can guide a choice of counterfactual, but cannot
establish what would have happened after the event. Pretrend tests may
have little power. Choosing an extrapolation because a specification
passes such a test also affects subsequent inference (Roth 2022).
Specify the reference period, transition window, counterfactual
assumption, and sensitivity range in advance.

\hypertarget{cohort-specific-estimation-and-aggregation}{%
\subsubsection{4.2. Cohort-specific estimation and
aggregation}\label{cohort-specific-estimation-and-aggregation}}

With staggered events, the cohorts contributing to an estimate may
change across event times. Start by estimating effects separately for
each cohort and event time, following the approach to group-time effects
in Callaway and Sant'Anna (2021). Within each cohort, estimate
conditional reduced forms and first stages for every institution or
other event unit. Form Wald ratios in covariate strata with adequate
support, standardize them, and then average across event units.

Taking ratios before averaging preserves the chosen weights. Pooling the
reduced forms and first stages first gives more weight to components
with larger first stages and generally changes the estimand. Use
prespecified weights to aggregate the cohort-specific effects. Keeping
the same cohorts and weights at each horizon makes comparisons across
event time easier to interpret.

Each cohort-period estimate still needs adequate first-stage and
covariate support, adequate outcome observation, and credible
assumptions relating its complier population to the others. If support
disappears, report the share of the reference distribution that is no
longer represented. Either justify extrapolating to that part of the
population or state that the estimand now covers only the supported set.

Stacking can place the same original observation in several cohort
comparisons. Those copies carry the same sampling information, so retain
the original cluster identifiers and rebuild the stacks within each
bootstrap resample. Assess calendar-time shocks when they change the
exposure contrast, compliance behavior, which outcomes are observed, or
instrument validity.

\hypertarget{estimation-inference-and-reporting}{%
\subsection{5. Estimation, inference, and
reporting}\label{estimation-inference-and-reporting}}

\hypertarget{estimation}{%
\subsubsection{5.1. Estimation}\label{estimation}}

Before looking at outcomes, fix the event date and transition window,
and define the target population, instrument, and exposure contrast.
Specify the covariates, support rule, and reference distribution. The
plan should explain how outcomes are observed, why the complier
populations can be compared, and how the counterfactual post-event LATE
will be constructed. For a staggered design, also fix the weights within
and across cohorts.

Equation (2) determines the calculation: estimate conditional reduced
forms and first stages, form Wald ratios where both periods have
adequate support, and average them using the reference weights. If
implementing this calculation through regression software, reproduce it
directly on the same sample with the same weights to check that the
software estimates the intended quantity.

Construct a generated instrument from predetermined information that the
event cannot affect. Fix the training population and dates, construction
rule, sign, cutoff, and observations or clusters to be omitted. Specify
how to handle inputs outside the prespecified support. If data from the
main analysis help estimate the score, rebuild it along with the rest of
the estimator during resampling. Leaving out the focal observation or
cluster prevents it from contributing mechanically to its own
instrument. Independence and exclusion still require an institutional
argument.

\Needspace{0.29\textheight}

\hypertarget{a-numerical-example}{%
\subsubsection{5.2. A numerical example}\label{a-numerical-example}}

Table 2 works through the calculation for two covariate groups with
reference weights of 0.60 and 0.40.

\textbf{Table 2. From means by instrument value to the standardized
contrast}

\begin{longtable}[]{@{}
  >{\raggedright\arraybackslash}p{(\columnwidth - 10\tabcolsep) * \real{0.1304}}
  >{\raggedleft\arraybackslash}p{(\columnwidth - 10\tabcolsep) * \real{0.1739}}
  >{\raggedleft\arraybackslash}p{(\columnwidth - 10\tabcolsep) * \real{0.1739}}
  >{\raggedleft\arraybackslash}p{(\columnwidth - 10\tabcolsep) * \real{0.1739}}
  >{\raggedleft\arraybackslash}p{(\columnwidth - 10\tabcolsep) * \real{0.1739}}
  >{\raggedleft\arraybackslash}p{(\columnwidth - 10\tabcolsep) * \real{0.1739}}@{}}
\toprule\noalign{}
\begin{minipage}[b]{\linewidth}\raggedright
Period
\end{minipage} & \begin{minipage}[b]{\linewidth}\raggedleft
Group
\end{minipage} & \begin{minipage}[b]{\linewidth}\raggedleft
Reference weight
\end{minipage} & \begin{minipage}[b]{\linewidth}\raggedleft
Reduced form
\end{minipage} & \begin{minipage}[b]{\linewidth}\raggedleft
First stage
\end{minipage} & \begin{minipage}[b]{\linewidth}\raggedleft
Wald ratio
\end{minipage} \\
\midrule\noalign{}
\endhead
\bottomrule\noalign{}
\endlastfoot
Pre & A & 0.60 & 0.36 & 0.20 & 1.80 \\
Pre & B & 0.40 & 0.24 & 0.16 & 1.50 \\
Post & A & 0.60 & 0.58 & 0.20 & 2.90 \\
Post & B & 0.40 & 0.40 & 0.16 & 2.50 \\
\end{longtable}

The pre-event standardized LATE is \(0.60(1.80)+0.40(1.50)=1.68\); after
the event, it is \(0.60(2.90)+0.40(2.50)=2.78\). The difference is 1.10.
Under the first three rows of Table 1, this is the standardized change
in period-specific LATEs. It is an event effect if the final row also
holds and the counterfactual post-event LATE equals the pre-event LATE.

\hypertarget{inference}{%
\subsubsection{5.3. Inference}\label{inference}}

The components of the final event estimate may share observations or
depend on the same reference weights, cohort stacks, or generated
instrument. Inference must account for their joint uncertainty,
including estimation of the counterfactual. When identification is
strong, one option is a cluster bootstrap that resamples the original
independent units and re-estimates the full contrast, including the
reference distribution, cohort stacks, and generated scores. Weak first
stages make the ratios unstable and require a different approach.

With weak first stages, inference must remain valid for the final event
estimand despite weak identification. Section E.3 gives a conservative
confidence-set procedure for one pre-event period, one post-event
period, fixed finite strata, fixed external weights, and a
counterfactual post-event LATE equal to the pre-event LATE. It first
constructs simultaneous intervals for the reduced forms and first
stages, then finds the event contrasts consistent with those intervals.
The set becomes unbounded when the data cannot exclude a zero
denominator.

Estimated weights, support selected from the data, and staggered
aggregation fall outside this procedure. They require inference for the
full estimator being used. Without such a procedure, report the reduced
forms and first stages and label conventional intervals as
approximations that require strong identification.

Choose clusters to reflect assignment and dependence, which may occur at
the ambulance-company, market, or hospital level, or at more than one of
these levels. Clustering only by patient may be insufficient. With few
clusters, a wild cluster bootstrap may improve inference for reduced
forms and first stages when its conditions hold (Cameron, Gelbach, and
Miller 2008). It does not resolve weak identification when applied to a
conventional ratio statistic.

\hypertarget{reporting}{%
\subsubsection{5.4. Reporting}\label{reporting}}

Report enough information for a reader to reconstruct each event
estimate: outcome and exposure means by instrument value, reduced forms,
first stages, reference weights, and standardized period-specific
ratios. Give observation rates and support losses separately, including
the share of the reference distribution retained. Identify the estimates
used to construct the counterfactual LATE. A staggered design should
also report cohort weights and clustering. For a generated instrument,
describe the training sample and state whether inference repeats the
steps used to construct it.

Use the interpretation supported by Table 1 when describing the result.
With the full identification argument, the estimate is a ``change in the
standardized LATE caused by the event for units who would comply in all
three states.'' If the evidence supports only weaker claims about
complier comparability or the counterfactual LATE, call it a
``standardized change in period-specific LATEs.''

\hypertarget{conclusion}{%
\subsection{6. Conclusion}\label{conclusion}}

In the hospital example, a change in the ambulance IV estimate can
reflect a change in care, a change in the patients whose destination
responds to ambulance assignment, or a change in which outcomes are
observed. The LATE may also have changed without acquisition.
Attributing the standardized IV difference to acquisition requires the
full set of conditions in Table 1 and a counterfactual LATE held at its
pre-event value. When the evidence supports a weaker interpretation,
report the period-specific LATEs and use sensitivity analysis to show
how assumptions about the missing counterfactual and the complier
populations affect the event estimate.

\vspace{2.2em}

\noindent

\rule{0.34\linewidth}{0.4pt}

\vspace{1.0em}

\textbf{Suggested citation}

Burla, Sriteja. 2026. \emph{Instrumental Variables Before and After an
Event: Identification with Endogenous Exposure}. Version 1.0, August
2026.

\newpage

\hypertarget{technical-supplement}{%
\section{Technical Supplement}\label{technical-supplement}}

\hypertarget{a.-setup-and-within-period-iv}{%
\subsection{A. Setup and within-period
IV}\label{a.-setup-and-within-period-iv}}

Consider one institution, or another unit at which the event occurs. Let
\(\mathcal U\) denote its target population. Membership in
\(\mathcal U\) is determined before the instrument and event. The index
\(a\in\{0,E\}\) records whether the event has occurred by the post-event
date: at \(t=1\), \(a=E\) is observed and \(a=0\) is counterfactual.
Under no anticipation (Section B.2), the two paths coincide at \(t=0\),
so the observed states are \(a_0=0\) and \(a_1=E\).

For individual \(i\) in period \(t\), let \(D_{it}(a,z)\in\{0,1\}\)
denote potential exposure under binary instrument value \(z\) and
\(Y_{it}(a,d)\) the potential outcome under exposure \(d\). This
notation imposes exclusion: holding exposure fixed, the instrument has
no further effect on the outcome. Exclusion is part of the IV
assumptions in each observed period. Extending it to the counterfactual
post-event state requires a separate assumption.

Keep the instrument coding fixed across periods, with \(z=1\) denoting
the value intended to increase exposure; monotonicity then requires
\(D_{it}(a,1)\geq D_{it}(a,0)\). Define the complier set and exposure
effect by

\[
\mathcal C_t(a)=\{i\in\mathcal U:D_{it}(a,1)>D_{it}(a,0)\},
\qquad
\tau_{it}(a)=Y_{it}(a,1)-Y_{it}(a,0).
\tag{S.1}
\]

Under the usual binary-IV conditions, Equation (1) identifies the
conditional mean effect of exposure among that period's observed
compliers (Imbens and Angrist 1994). Section 2.1 discusses how to assess
these conditions and report estimates for each period.

The indicator \(O_{it}\) records whether an eligible unit's outcome is
observed. Conditional on predetermined covariates, the benchmark assumes
that neither the instrument nor exposure affects outcome observation. If
either affects survival, attrition, or missingness, the analysis needs a
separate argument to address selection.

\hypertarget{b.-standardization-and-event-identification}{%
\subsection{B. Standardization and event
identification}\label{b.-standardization-and-event-identification}}

\hypertarget{b.1.-standardization-and-outcome-observation}{%
\subsubsection{B.1. Standardization and outcome
observation}\label{b.1.-standardization-and-outcome-observation}}

Let \(\mathcal X^\star\) contain the covariate values supported in both
observed periods, with reference distribution \(F_X^\star\) as defined
in Section 2.2.

To extend the result to all compliers in the target population, assume
that their conditional mean effect equals that of compliers whose
outcomes are observed:

\[
E[\tau_{it}(a_t)\mid i\in\mathcal C_t(a_t),X_i=x,O_{it}=1]
=E[\tau_{it}(a_t)\mid i\in\mathcal C_t(a_t),X_i=x].
\tag{S.2}
\]

Under Equation (S.2), the standardized Wald ratio averages the
conditional effects of exposure among all compliers in the target
population, using \(F_X^\star\):

\[
\beta_t^\star
=E_{F_X^\star}\!\left[
E[\tau_{it}(a_t)\mid i\in\mathcal C_t(a_t),X_i]
\right].
\tag{S.3}
\]

Equation (S.2) restricts only the conditional mean. Equating the joint
distributions of potential exposures and outcomes for observed and
target-population compliers would require a stronger assumption. Both
versions require a positive probability of observing outcomes wherever
\(F_X^\star\) assigns weight.

\hypertarget{b.2.-a-common-complier-population-and-the-post-event-counterfactual}{%
\subsubsection{B.2. A common complier population and the post-event
counterfactual}\label{b.2.-a-common-complier-population-and-the-post-event-counterfactual}}

Let \(\mathcal C^\cap\) be the intersection of the complier sets before
the event and under the two possible post-event states:

\[
\mathcal C^\cap
=\mathcal C_0(0)\cap\mathcal C_1(0)\cap\mathcal C_1(E).
\tag{S.4}
\]

The main result assumes that all three sets equal \(\mathcal C^\cap\) at
every value of \(X\) with positive reference weight. For these
compliers, define

\[
L_t^{a,\cap,\star}
=E_{F_X^\star}\!\left[
E[\tau_{it}(a)\mid i\in\mathcal C^\cap,X_i]
\right].
\tag{S.5}
\]

No anticipation means that potential exposure and outcomes before the
event cannot depend on the future event path:

\[
D_{i0}(E,z)=D_{i0}(0,z),
\qquad
Y_{i0}(E,d)=Y_{i0}(0,d)
\]

for every relevant \(z\) and \(d\). These equalities allow the observed
pre-event quantities in Equations (S.4)--(S.6) to be written under the
no-event state. To specify the missing \(L_1^{0,\cap,\star}\),
researchers also need a rule that uses the comparable pre-event LATE
history. The simplest rule assumes
\(L_1^{0,\cap,\star}=L_0^{0,\cap,\star}\).

Exclusion must hold in the counterfactual post-event state as well.
Because this state is unobserved, the justification must come from
institutional evidence.

\textbf{Identification result.} Suppose period-specific IV validity and
Equation (S.2) hold. If the three complier sets are equal, there is no
anticipation, exclusion extends to the counterfactual post-event state,
and the counterfactual post-event LATE equals the pre-event LATE, then

\[
\beta_1^\star-\beta_0^\star
=L_1^{E,\cap,\star}-L_0^{0,\cap,\star}
=L_1^{E,\cap,\star}-L_1^{0,\cap,\star}
=\theta^{\cap,\star}.
\tag{S.6}
\]

A different assumption about the counterfactual LATE replaces
\(\beta_0^\star\) with a prespecified value constructed from the
pre-event LATE history.

\hypertarget{b.3.-allowing-the-complier-population-to-change}{%
\subsubsection{B.3. Allowing the complier population to
change}\label{b.3.-allowing-the-complier-population-to-change}}

Identification can allow the three complier sets to contain different
units if the relevant mean effects are equal. To state this weaker
condition, write, for any group \(G\),

\[
M_{t,a}^{G}(x)
=E[\tau_{it}(a)\mid i\in G,X_i=x].
\]

For this identification argument, impose the following conditional mean
equalities:

\[
\begin{aligned}
M_{0,0}^{\mathcal C_0(0)}(x)
&=M_{0,0}^{\mathcal C_1(0)}(x),\\
M_{1,0}^{\mathcal C_1(0)}(x)
&=M_{1,0}^{\mathcal C^\cap}(x),\\
M_{1,E}^{\mathcal C_1(E)}(x)
&=M_{1,E}^{\mathcal C^\cap}(x).
\end{aligned}
\tag{S.7}
\]

The first equality says that pre-event compliers have the same pre-event
mean effect as those who would comply at the post-event date without the
event. The counterfactual assumption then specifies how the mean effect
for this latter population would evolve. Assuming it remains at its
pre-event value gives
\(M_{1,0}^{\mathcal C_1(0)}(x)=M_{0,0}^{\mathcal C_1(0)}(x)\).

The second equality says that the counterfactual post-event mean effect
also holds for \(\mathcal C^\cap\). The third imposes the corresponding
equality for the post-event mean under the event. These restrictions
allow the complier populations to differ in membership.

All three restrictions apply wherever \(F_X^\star\) assigns weight, and
\(\mathcal C^\cap\) must have positive probability there. When group
membership differs, equal conditional mean effects need a substantive
justification.

\hypertarget{c.-decomposition-and-sensitivity}{%
\subsection{C. Decomposition and
sensitivity}\label{c.-decomposition-and-sensitivity}}

\hypertarget{c.1.-decomposition}{%
\subsubsection{C.1. Decomposition}\label{c.1.-decomposition}}

Let \(m_t^O\) denote the conditional effects among observed compliers in
period \(t\), averaged over the reference distribution. Under
period-specific IV validity, \(m_t^O=\beta_t^\star\). Let \(m_t^U\)
denote the corresponding average for all period-\(t\) compliers in the
target population. The three discrepancies in Equation (5) are

\[
\begin{aligned}
\Delta^{0,\star}
&=L_1^{0,\cap,\star}-L_0^{0,\cap,\star},\\
B^{C,\star}
&=(m_1^U-L_1^{E,\cap,\star})-(m_0^U-L_0^{0,\cap,\star}),\\
B^{O,\star}
&=(m_1^O-m_1^U)-(m_0^O-m_0^U).
\end{aligned}
\tag{S.8}
\]

Adding and subtracting these effects gives Equation (5). A constant
counterfactual LATE for the common complier population sets
\(\Delta^{0,\star}=0\). Equality of the three complier sets sets
\(B^{C,\star}=0\), and Equation (S.2) sets \(B^{O,\star}=0\).

Under Section B.3's conditional mean equalities and its assumption that
the no-event mean effect for \(\mathcal C_1(0)\) is unchanged between
periods, \(\Delta^{0,\star}+B^{C,\star}=0\). Neither term need be zero
separately.

\hypertarget{c.2.-sensitivity-region}{%
\subsubsection{C.2. Sensitivity region}\label{c.2.-sensitivity-region}}

Let \(\mathcal H\) contain the prespecified plausible triples
\((\delta,b_C,b_O)\) for these discrepancies. The corresponding set of
event effects is

\[
\Theta(\mathcal H)
=\left\{
\Delta^{IV,\star}-\delta-b_C-b_O:
(\delta,b_C,b_O)\in\mathcal H
\right\}.
\tag{S.9}
\]

Equation (6) uses symmetric rectangular bounds. More generally,
\(\mathcal H\) can allow asymmetric bounds, impose signs, or restrict
how the discrepancies vary together. Historical IV estimates can inform
the counterfactual LATE component when they use the same instrument,
exposure, target population, reference distribution, and assumptions
about complier comparability.

\hypertarget{d.-multiple-periods-and-staggered-events}{%
\subsection{D. Multiple periods and staggered
events}\label{d.-multiple-periods-and-staggered-events}}

With multiple periods, apply the two-period argument separately within
each cohort at each event time, then aggregate the estimates.

Let \(c\) index event cohorts, \(h\) institutions or other event units,
and \(k\) event time. Write the conditional reduced form and first stage
for each component as \(\rho_{hck}(x)\) and \(\pi_{hck}(x)\). Using
reference distribution \(F_{X,hc}^\star\), its standardized LATE in the
observed period is

\[
\beta_{hck}^\star
=E_{F_{X,hc}^\star}\!\left[
\frac{\rho_{hck}(X)}{\pi_{hck}(X)}
\right].
\tag{S.10}
\]

For a reported post-event horizon, define
\(\mathcal C_{hck}^\cap=\mathcal C_{hck}(0)\cap\mathcal C_{hck}(E)\).
Construct the prespecified counterfactual LATE, \(g_{hck}\), from
comparable pre-event LATEs for the same complier population and
reference distribution. Period-specific IV validity and Equation (S.2)
identify the conditional effects among target-population compliers in
the observed state, averaged over \(F_{X,hc}^\star\). The required
conditional mean equalities allow both the observed and counterfactual
effects to apply to \(\mathcal C_{hck}^\cap\). If there is no
anticipation, exclusion holds in the counterfactual state, and
\(g_{hck}\) is correctly specified,

\[
\beta_{hck}^\star-g_{hck}
=\theta_{hck}^{\cap,\star}.
\tag{S.11}
\]

The complier population in (S.11) can differ across values of \(k\). To
interpret the full event-study path for a fixed complier population, use
the intersection across every reported horizon. Support must hold in the
reported period and in each period used to construct \(F_{X,hc}^\star\)
or \(g_{hck}\). It therefore has to be assessed separately for each
horizon.

Let \(\mathcal G_k\) denote the cohorts that contribute at horizon
\(k\), and \(\mathcal H_c\) the event units in cohort \(c\). The cohort
weights are nonnegative and satisfy
\(\sum_{c\in\mathcal G_k}\omega_{c\mid k}=1\); within a cohort,
\(\sum_{h\in\mathcal H_c}\omega_{hc\mid k}=1\). The aggregate is

\[
\theta_k
=\sum_{c\in\mathcal G_k}\omega_{c\mid k}
\sum_{h\in\mathcal H_c}\omega_{hc\mid k}
\left\{
\beta_{hck}^\star
-g_{hck}
\right\}.
\tag{S.12}
\]

Pooling reduced forms and first stages before taking their ratio changes
the weights. For generic components \(j\) with prespecified weights
\(a_j\),

\[
\frac{\sum_j a_j\rho_j}{\sum_j a_j\pi_j}
=\sum_j
\frac{a_j\pi_j}{\sum_\ell a_\ell\pi_\ell}
\frac{\rho_j}{\pi_j}.
\tag{S.13}
\]

The pooled ratio therefore weights each component according to its first
stage, changing the weights in (S.12). The target also changes across
horizons unless \(\mathcal G_k\) and the cohort weights remain fixed
across \(k\). Dropping a component that lacks support changes the target
as well. Report which cohorts contribute, their weights, and the share
of the reference distribution retained.

\hypertarget{e.-estimation-and-inference}{%
\subsection{E. Estimation and
inference}\label{e.-estimation-and-inference}}

\hypertarget{e.1.-finite-strata-estimator}{%
\subsubsection{E.1. Finite-strata
estimator}\label{e.1.-finite-strata-estimator}}

Compute a Wald ratio for each period in each stratum with adequate
support, then average the ratios using the chosen reference weights.

For finite covariate strata \(x\in\mathcal X^\star\), let
\(\widehat p_x^\star\) be the reference weight and \(\widehat g_1\) the
estimated counterfactual post-event LATE without the event. This gives
the plug-in estimator

\[
\widehat\beta_t^\star
=\sum_{x\in\mathcal X^\star}
\widehat p_x^\star
\frac{\widehat\rho_t(x)}{\widehat\pi_t(x)},
\qquad
\widehat\theta=\widehat\beta_1^\star-\widehat g_1.
\tag{S.14}
\]

In the simplest two-period case, \(\widehat g_1=\widehat\beta_0^\star\).
Other counterfactual assumptions may use several pre-event estimates.
Apply the support rule and weights specified for the target population.

\Needspace{0.14\textheight}

The default reference is the pre-event complier distribution. Let
\(\widehat f_0(x)\) denote the empirical pre-event covariate mass over
the supported strata. The estimated reference weights are

\[
\widehat p_x^\star
=\frac{\widehat\pi_0(x)\widehat f_0(x)}
{\sum_v\widehat\pi_0(v)\widehat f_0(v)}.
\tag{S.15}
\]

For the pre-event component, substituting (S.15) cancels the first
stages within each stratum:

\[
\widehat\beta_0^\star
=\frac{\sum_x\widehat f_0(x)\widehat\rho_0(x)}
{\sum_v\widehat f_0(v)\widehat\pi_0(v)}.
\]

The post-event component still contains the ratios
\(\widehat\rho_1(x)/\widehat\pi_1(x)\) in the supported strata; its
first stages do not cancel. When the reference weights come from the
analysis data, inference under strong identification should re-estimate
them. Section E.3 considers fixed weights obtained from an external
source.

\hypertarget{e.2.-strong-identification-inference}{%
\subsubsection{E.2. Strong-identification
inference}\label{e.2.-strong-identification-inference}}

Resample the original independent units and reconstruct every component
determined by the data: the reference distribution, generated
instrument, cohort stacks, counterfactual LATE, and aggregation weights.
This preserves covariance among the components of \(\widehat\theta\).
Prespecified support can remain fixed. If support is selected from the
data, inference must account for that selection, and researchers must
state which population the estimates represent.

Calculate the final contrast in each resample and state how its
bootstrap distribution is used to construct a confidence interval. For
an event-time path, reconstruct all horizons within each resample. Use
the distribution at each horizon for pointwise intervals and their joint
distribution for a simultaneous band. Analytic variance calculations
must also retain covariance among the components. A sandwich variance
requires stacked scores for the estimator that averages the conditional
ratios. Pooling the reduced forms and first stages before taking their
ratio changes both the estimand and its influence function.

\hypertarget{e.3.-weak-instrument-inference}{%
\subsubsection{E.3. Weak-instrument
inference}\label{e.3.-weak-instrument-inference}}

The following construction can be implemented in the limited case of one
pre-event period, one post-event period, fixed finite strata, fixed
external reference weights, and a counterfactual post-event LATE equal
to the pre-event LATE.

Suppose there are \(K\) positively weighted strata and \(m=4K\)
reduced-form and first-stage moments. For each moment \(q_j\), construct
the interval

\[
\mathcal I_j=
\left[
\widehat q_j-z_{1-\alpha/(2m)}
\widehat{\operatorname{se}}(\widehat q_j),
\quad
\widehat q_j+z_{1-\alpha/(2m)}
\widehat{\operatorname{se}}(\widehat q_j)
\right].
\tag{S.16}
\]

If the marginal normal approximations are valid under the stated
sampling and clustering design, Bonferroni's inequality gives this
rectangle joint coverage of at least \(1-\alpha\). The standard errors
must account for the application's dependence structure. The normal
critical value in (S.16) does not correct inference with few clusters.

\Needspace{0.21\textheight}

Within each period and stratum, form a set of ratios from the
reduced-form and first-stage intervals. If the first-stage interval
excludes zero, calculate the four quotients of their endpoints; the
smallest and largest quotients bound the ratio interval. If the
first-stage interval includes zero, use the whole real line. Denote the
resulting set by \(\mathcal B_t(x)\). The confidence set for the event
contrast is

\[
\mathrm{CS}_{1-\alpha}
=\sum_x p_x^\star
\{\mathcal B_1(x)-\mathcal B_0(x)\}.
\tag{S.17}
\]

This construction is conservative because it ignores covariance that
could narrow the joint set. Weak first stages widen the set and can make
it unbounded. The worked files supplied with this guide use fixed
external weights and give a point estimate of 1.00, with a conservative
95 percent set of approximately \([-1.37,3.50]\). Despite the point
estimate, this wide set leaves substantial uncertainty about the event
contrast.

First-stage intervals also help assess the design. A negative point
estimate warrants attention, although it is not a test of monotonicity.
An interval spanning zero provides weak and inconclusive evidence. If
the entire interval lies at or below zero, the data contradict the
assumed positive first stage at the reported confidence level. A causal
interpretation then requires recoding the instrument or redefining and
defending the design.

The companion implementation reproduces Equations (S.16)--(S.17),
including the worked input and reported result. It does not cover
estimated reference weights, strata or support chosen from the data,
flexible covariates, nonlinear extrapolation of the counterfactual LATE,
overlapping cohort stacks, or estimated cohort weights. These cases
require identification-robust inference for the estimator used. If no
justified procedure is available, report the reduced forms and first
stages and label conventional intervals as strong-identification
approximations. Anderson--Rubin procedures remain useful for suitable
component equations. To obtain valid coverage for a difference between
dependent components, however, the procedure must handle them jointly
(Anderson and Rubin 1949; Andrews 2022).

\hypertarget{e.4.-dependence-and-stacked-data}{%
\subsubsection{E.4. Dependence and stacked
data}\label{e.4.-dependence-and-stacked-data}}

Variance estimation must account for assignment, sampling, event timing,
and instrument construction. Repeated observations on the same person
and copies of an observation in different cohort stacks do not provide
independent information. Keep the original cluster identifiers and
rebuild the stacks in each resample. With few independent clusters, use
a procedure justified for the sampling structure.

A cluster bootstrap of a conventional ratio statistic still requires
strong identification. The projection in Section E.3 allows weak
instruments but relies on credible intervals for the reduced-form and
first-stage moments. It cannot correct size distortions caused by few
clusters. The two problems require separate treatment.

\Needspace{0.22\textheight}

\hypertarget{e.5.-generated-instruments}{%
\subsubsection{E.5. Generated
instruments}\label{e.5.-generated-instruments}}

Suppose the scoring rule was estimated in an external training sample.
The analysis can condition on that sample and the estimated rule and
treat the score as fixed. Unconditional inference must also account for
uncertainty in the estimated scoring rule. Keep the rule's input
definition, scale, sign, and any cutoff unchanged across periods: these
features define the encouragement intervention and therefore the
complier population.

If the score also uses data from the main analysis, resample the
original independent units and rebuild the training sample, score,
cutoff, and assignment rule. Choose which observations to leave out
according to the source of dependence. Omitting only the focal
observation may be insufficient when patients share an ambulance
company, market, event unit, or training cluster. Leaving out the
relevant observation or cluster prevents it from mechanically
contributing to its own instrument. Independence and exclusion still
require an institutional argument.

\Needspace{0.16\textheight}

New assignment sources or input patterns may appear after the event. A
prespecified rule can classify them as unsupported, map them using
supported predetermined features, or examine them separately with an
instrument re-estimated after the event. Retraining the primary score on
post-event exposure choices can make the instrument itself respond to
the event. For each period and cohort, report score support, shares by
instrument value, first stages, and the share of the reference
distribution lacking support.

\begingroup
\small
\setlength{\parskip}{0.28em}

\hypertarget{references}{%
\subsection{References}\label{references}}

Anderson, T. W., and Herman Rubin. 1949. ``Estimation of the Parameters
of a Single Equation in a Complete System of Stochastic Equations.''
\emph{Annals of Mathematical Statistics} 20 (1): 46--63.
\url{https://doi.org/10.1214/aoms/1177730090}.

Andrews, Donald W. K. 2022. ``Identification-Robust Subvector
Inference.'' Cowles Foundation Discussion Paper 2105. Originally issued
2017. \url{https://cowles.yale.edu/node/139273}.

Callaway, Brantly, and Pedro H. C. Sant'Anna. 2021.
``Difference-in-Differences with Multiple Time Periods.'' \emph{Journal
of Econometrics} 225 (2): 200--230.
\url{https://doi.org/10.1016/j.jeconom.2020.12.001}.

Cameron, A. Colin, Jonah B. Gelbach, and Douglas L. Miller. 2008.
``Bootstrap-Based Improvements for Inference with Clustered Errors.''
\emph{Review of Economics and Statistics} 90 (3): 414--427.
\url{https://doi.org/10.1162/rest.90.3.414}.

Doyle, Joseph J., Jr., John A. Graves, Jonathan Gruber, and Samuel A.
Kleiner. 2015. ``Measuring Returns to Hospital Care: Evidence from
Ambulance Referral Patterns.'' \emph{Journal of Political Economy} 123
(1): 170--214. \url{https://doi.org/10.1086/677756}.

de Chaisemartin, Clément, and Xavier D'Haultfœuille. 2018. ``Fuzzy
Differences-in-Differences.'' \emph{Review of Economic Studies} 85 (2):
999--1028. \url{https://doi.org/10.1093/restud/rdx049}.

Imbens, Guido W., and Joshua D. Angrist. 1994. ``Identification and
Estimation of Local Average Treatment Effects.'' \emph{Econometrica} 62
(2): 467--475. \url{https://doi.org/10.2307/2951620}.

Konetzka, R. Tamara, Fan Yang, and Rachel M. Werner. 2019. ``Use of
Instrumental Variables for Endogenous Treatment at the Provider Level.''
\emph{Health Economics} 28 (5): 710--716.
\url{https://doi.org/10.1002/hec.3861}.

Miyaji, Sho. 2024. ``Instrumented Difference-in-Differences with
Heterogeneous Treatment Effects.'' RIEB Discussion Paper Series 2024-22.
Kobe University.
\url{https://www.rieb.kobe-u.ac.jp/academic/ra/dp/English/dp2024-22.html}.

Rambachan, Ashesh, and Jonathan Roth. 2023. ``A More Credible Approach
to Parallel Trends.'' \emph{Review of Economic Studies} 90 (5):
2555--2591. \url{https://doi.org/10.1093/restud/rdad018}.

Roth, Jonathan. 2022. ``Pretest with Caution: Event-Study Estimates
after Testing for Parallel Trends.'' \emph{American Economic Review:
Insights} 4 (3): 305--322. \url{https://doi.org/10.1257/aeri.20210236}.

\endgroup

\end{document}
